% SPDX-License-Identifier: CC-BY-4.0
\documentclass[11pt,letterpaper]{article}

\usepackage[T1]{fontenc}
\usepackage{amsmath,amssymb}
\usepackage{tikz}
\usepackage{pgfplots}
\pgfplotsset{compat=1.18}
\usepackage{algorithm}
\usepackage{algpseudocode}
\usepackage{booktabs}
\usepackage{enumitem}
\usepackage[hidelinks]{hyperref}

\newcommand{\N}{\mathbb{N}}
\newcommand{\Z}{\mathbb{Z}}
\newcommand{\Q}{\mathbb{Q}}
\newcommand{\R}{\mathbb{R}}
\newcommand{\C}{\mathbb{C}}
\newcommand{\Map}[3]{#1 \colon #2 \to #3}
\newcommand{\Mod}[1]{#1\text{-}\mathrm{Mod}}
\newcommand{\GroupHom}{\mathrm{Hom}_{\mathbf{Grp}}}
\newcommand{\RingHom}{\mathrm{Hom}_{\mathbf{Ring}}}

\newtheorem{theorem}{Theorem}
\newtheorem{lemma}[theorem]{Lemma}
\newtheorem{corollary}[theorem]{Corollary}
\newtheorem{definition}[theorem]{Definition}
\newtheorem{remark}[theorem]{Remark}
\newenvironment{proof}{\noindent\textit{Proof.}}{\hfill\(\square\)}

\title{Generalized Collatz Library}
\author{}
\date{\today}

\begin{document}
\maketitle
\section*{Exact Gaussian Division and Homogeneous Escape}

\section{Ring domain and exact division}

Work in \(R=\Z[i]\), with \(N(x+iy)=x^2+y^2\). For \(d\ne0\), define the division-precedence raw map
\[
T(z)=
\begin{cases}
z/d,&d\mid z\text{ in }R,\\
qz+c,&d\nmid z\text{ in }R.
\end{cases}
\]
The quotient is used only on the exact division branch. This is not divisibility in the field \(\C\).

% The adjugate test retains both integral coordinates of the quotient.
\begin{theorem}[Exact Gaussian division]
For \(d=s+it\ne0\) and \(z=x+iy\),
\[
d\mid z\quad\Longleftrightarrow\quad
(s^2+t^2)\mid sx+ty\ \text{ and }\ (s^2+t^2)\mid sy-tx.
\]
\end{theorem}
\begin{proof}
Multiplication by \(s+it\) has matrix
\[
\begin{pmatrix}s&-t\\t&s\end{pmatrix},
\qquad \det=s^2+t^2>0.
\]
Inverting it gives
\[
\frac{z}{d}=\frac{sx+ty}{s^2+t^2}
+i\,\frac{sy-tx}{s^2+t^2}.
\]
If \(z=dw\) with \(w\in R\), the two displayed coefficients are integers. Conversely, the two divisibilities make those coefficients integers; multiplication of the resulting Gaussian integer by \(d\) gives \(z\). This proves necessity and sufficiency, and is checked by \texttt{DvdCoordinates}.
\end{proof}

\section{The well-founded division prefix}

% A positive integral squared norm decreases at every unresolved prefix step.
\begin{theorem}[Division prefix and integral budget]
Suppose \(N(d)\ge2\) and \(z\ne0\). For arbitrary \(q,c\in R\), there are \(r\in\N\) and \(t\in R\) satisfying
\[
d\nmid t,\qquad T^r(z)=t,\qquad z=d^rt,\qquad 2^r\le N(z).
\]
Before raw time \(r\), every transition is division.
\end{theorem}
\begin{proof}
Norm is multiplicative:
\[
N(uv)=N(u)N(v).
\]
This follows either by multiplying \(u\bar u\,v\bar v\), or by expanding both real and imaginary coordinates. If \(d\mid w\) with \(w\ne0\), exact division gives \(w=d(w/d)\). Its quotient cannot be zero. Consequently
\[
N(w)=N(d)N(w/d)\ge2N(w/d)>N(w/d)>0.
\]
Strong induction on the positive integer \(N(w)\) proves termination. At a nondivisible \(w\), use \(r=0,t=w\). At a divisible \(w\), invoke the induction hypothesis on \(w/d\), prepend its one division tick, and multiply its factorization by \(d\). This reconstructs both \(T^r(z)=t\) and \(z=d^rt\).

The final \(t\) is nonzero, since \(d\mid0\). Therefore
\[
N(z)=N(d)^rN(t)\ge2^r.
\]
The Lean proofs are \texttt{DivisionPrefix} and \texttt{DivisionPrefixBudget}. This is an exact finite prefix budget, not a floating logarithmic estimate or a claim of polynomial running time in bit length.
\end{proof}

\section{A sufficient avoidance invariant}

% One common Bezout witness rules out every future division in the same orbit.
\begin{theorem}[Homogeneous escape]
Assume \(c=0\), \(N(q),N(d)\ge2\), and supply \(a,b\in R\) with \(aq+bd=1\). Every nonzero raw seed escapes in full sequence. With the division prefix above,
\[
T^{r+j}(z)=q^jt,\qquad
N(T^{r+j}(z))=N(q)^jN(t)\ge2^j.
\]
\end{theorem}
\begin{proof}
Suppose \(d\nmid u\). If \(d\mid qu\), then
\[
u=(aq+bd)u=a(qu)+b(du)
\]
would be divisible by \(d\), a contradiction. This single identity proves that multiplication by \(q\) preserves avoidance. Induction gives \(d\nmid q^jt\) for every \(j\). Hence every raw tail step is the homogeneous affine branch, proving the state formula by induction at the original clock. Multiplicativity gives the exact norm formula.

For any \(B\in\N\), choose \(K=r+B+1\). If \(k\ge K\), write \(k=r+j\) with \(j\ge B+1\). The elementary induction \(j<2^j\) gives
\[
B<j<2^j\le N(T^k(z)).
\]
Thus every sufficiently late state, not just a subsequence, exceeds the bound. This is \texttt{CoprimeHomogeneousEscape}. The required identity is not replaced by an assertion about common element divisors in an arbitrary quadratic order.
\end{proof}

\section{Equal norms do not imply a common base}

% Conjugate split primes separate by an explicit residue homomorphism.
\begin{corollary}[Conjugate split primes]
Let \(q=2+i\), \(d=2-i\), and \(c=0\). Then \(N(q)=N(d)=5\), but every nonzero raw seed escapes, and \(q^a\ne d^b\) for all positive integers \(a,b\).
\end{corollary}
\begin{proof}
Direct multiplication gives
\[
-q+(1+i)d=1.
\]
The preceding theorem applies. Zero is separately fixed by division precedence.

The map \(R\to\mathbb F_5\) sending \(i\mapsto2\) is a ring homomorphism because \(2^2+1=0\) modulo \(5\). It sends \(q\mapsto4\), a unit, and \(d\mapsto0\). Any positive power of \(d\) maps to zero, whereas every power of \(q\) maps to a unit. Their equality is impossible. Lean checks these facts as \texttt{EqualNormSplitEscape}, \texttt{SplitResidue}, and \texttt{EqualNormNoCommonPower}.
\end{proof}

\section{Reduction contract and scope}

The original obligation is the full tail of each nonzero homogeneous Gaussian orbit under the specified norm and common-witness hypotheses. The replacement removes a finite division prefix, proves one guard-avoidance invariant, and computes a geometric-power tail. Positive integral height is a well-founded decreasing resource on every unresolved prefix child; the budget \(2^r\le N(z)\) bounds its length. Recovery retains the same Gaussian seed, exact ring guard, and every original raw tick.

No additive-offset classification, noncoprime case, unit-divisor case, or universal integer Collatz convergence follows. Equal norms do not identify ideals or elements. A lattice census is an illustration, not the finite-to-infinite proof. The companion complex survey separately records the source conventions, cohomology calculations, and the missing Julia-to-arithmetic bridge.

\end{document}
